\chapter{ความรู้พื้นฐานทางทฤษฎีจำนวน}
\begin{definitionbox}
	ให้ $x$ เป็นจำนวนจริง ค่าสัมบูรณ์ของ $x$ เขียนเป็น $|x|$ และนิยามได้ดังนี้
	\[
	|x| =
	\begin{cases}
		x, & \text{ถ้า}\; x \ge 0, \\
		-x, & \text{ถ้า}\; x < 0
	\end{cases}
	\]    
\end{definitionbox}

\begin{examplebox}
	\begin{enumerate}
		\item $|7|=$
		\item $|-3|=$
		\item $|-(-2)|=$
		\item $|(-3)(-5)|=$
	\end{enumerate}  
\end{examplebox}

\newpage
\begin{theorembox}[ขั้นตอนวิธีการหาร (Division Algorithm)]
	ให้ $a$ และ $b$ เป็นจำนวนเต็มโดยที่ $a \neq 0$ แล้วจะมีจำนวนเต็ม $q$ และ $r$ เพียงคู่เดียวเท่านั้น ที่ทำให้ 
	\[b=a q+r\]
	โดยที่ $0 \leq r<|a|$ จะเรียก $q$ ว่า \textbf{ผลหาร} (quotient) และ $r$ ว่า \textbf{เศษ} (remainder)   
\end{theorembox}

\newpage
\begin{examplebox}
	ให้ $b=60$ และ $a=7$ จงหาจำนวนเต็ม $q$ และ $r$ ตามขั้นตอนวิธีการหาร ที่ทำให้\\
	\[b=aq+r\;\text{โดยที่ $0 \leq r<|a|$}\]
\end{examplebox}
\solution เลือก $q=8$ และ $r=4$ เห็นได้ว่าสอดคล้องกับ $b=aq+r$ โดยที่ $0 \leq r<|a|$ เนื่องจาก
\[60=7(8)+4, \quad 0 \leq 4<|7|\]

\begin{examplebox}
	ให้ $b=-39$ และ $a=6$ จงหาจำนวนเต็ม $q$ และ $r$ ตามขั้นตอนวิธีการหาร ที่ทำให้\\
	\[b=aq+r\;\text{โดยที่ $0 \leq r<|a|$}\]

\end{examplebox}
\solution เลือก $q=-7$ และ $r=3$ เห็นได้ว่าสอดคล้องกับ $b=aq+r$ โดยที่ $0 \leq r<|a|$ เนื่องจาก
\[
-39=6(-7)+3, \quad 0 \leq 3<|6|
\]

\begin{examplebox}
	ให้ $b=52, a=5$ จงหาจำนวนเต็ม $q$ และ $r$ ตามขั้นตอนวิธีการหาร ที่ทำให้\\
	\[b=aq+r\;\text{โดยที่ $0 \leq r<|a|$}\]  
\end{examplebox}
\solution
\vspace{2cm}  

\begin{examplebox}
	ให้ $b=-52, a=5$ จงหาค่า $q$ และ $r$ ตามขั้นตอนวิธีการหาร ที่ทำให้ \\
	\[b=aq+r\;\text{โดยที่ $0 \leq r<|a|$}\] 
\end{examplebox}
	\solution
\vspace{2cm}  

\newpage
\begin{examplebox}
	ให้ $b=52, a=-5$ จงหาจำนวนเต็ม $q$ และ $r$ ตามขั้นตอนวิธีการหาร ที่ทำให้ \\
	\[b=aq+r\;\text{โดยที่ $0 \leq r<|a|$}\]
\end{examplebox}
	\solution 
\vspace{2cm} 

\begin{examplebox}
	ให้ $b=-52, a=-5$ จงหาจำนวนเต็ม $q$ และ $r$ ตามขั้นตอนวิธีการหาร ที่ทำให้ \\
	\[b=aq+r\;\text{โดยที่ $0 \leq r<|a|$}\] 
\end{examplebox}
\solution
\vspace{2cm}  

\begin{examplebox}
	ให้ $b=5$ และ $a=3$ จงหาจำนวนเต็ม $q$ และ $r$ ตามขั้นตอนวิธีการหาร ที่ทำให้ \\
	\[b=aq+r\;\text{โดยที่ $0 \leq r<|a|$}\]
\end{examplebox}
\solution

\begin{examplebox}
	ให้ $b=59$ และ $a=-14$ จงหาจำนวนเต็ม $q$ และ $r$ ตามขั้นตอนวิธีการหาร ที่ทำให้ \\
	\[b=aq+r\;\text{โดยที่ $0 \leq r<|a|$}\] 
\end{examplebox}
	\solution
\vspace{9cm} 


\begin{examplebox}
	ให้ $b=-79$ และ $a=11$ จงหาจำนวนเต็ม $q$ และ $r$ ตามขั้นตอนวิธีการหาร ที่ทำให้ \\
	\[b=aq+r\;\text{โดยที่ $0 \leq r<|a|$}\]
\end{examplebox}
\solution


\newpage
\section{การหารลงตัว}
\begin{definitionbox}
	ให้ $a$ และ $b$ เป็นจำนวนเต็มโดยที่ $a \neq 0$
	\begin{itemize}
		\item เรียก $a$ ว่าเป็น \textbf{ตัวหาร} (divisor) หรือ \textbf{ตัวประกอบ} (factor) ตัวหนึ่งของ $b$ ก็ต่อเมื่อ มีจำนวนเต็ม $q$ ที่ทำให้ $b=a q$
		\item ในกรณีที่ $a$ เป็นตัวหารของ $b$ เราจะเรียกอีกแบบหนึ่งว่า $a$ หาร $b$ ลงตัว สามารถเขียนแทนได้ด้วยสัญลักษณ์ $a \mid b$
		\item ในกรณีที่ $a$ หาร $b$ ไม่ลงตัว สามารถเขียนแทนได้ด้วยสัญลักษณ์ $a \nmid b$
	\end{itemize}
\end{definitionbox}

\begin{examplebox}
	\begin{enumerate}
		\item $13 \mid 182$ เพราะมีจำนวนเต็ม $14$ ที่ทำให้ $182=(13)(14)$
		\item $-5 \mid 30$ เพราะมีจำนวนเต็ม $-6$ ที่ทำให้ $30=(-5)(-6)$
		\item $2 \nmid 5$ เพราะไม่มีจำนวนเต็ม $c$ ที่ทำให้ $5=2c$
	\end{enumerate}    
\end{examplebox}

\subsection{สมบัติการหารลงตัวที่สำคัญ}
\begin{theorembox}
	ให้ $a, b$ และ $c$ เป็นจำนวนเต็ม จะได้ว่า
	\begin{enumerate}
		\item $a \mid 0$
		\item $1\mid a$
		\item $a \mid a$
		\item $a \mid 1$ ก็ต่อเมื่อ $a= \pm 1$
		\item ถ้า $a \mid b$ และ $b \mid c$ แล้ว $a \mid c$
		\item ถ้า $a \mid b$ และ $c \mid d$ แล้ว $a c \mid b d$
		\item ถ้า $a \mid(b+c)$ และ $a \mid b$ แล้ว $a \mid c$
		\item $a \mid b$ และ $b \mid a$ ก็ต่อเมื่อ $a= \pm b$
		\item ถ้า $a \mid b$ และ $b \neq 0$ แล้ว $|a| \leq|b|$
		\item ถ้า $a \mid b$ และ $a \mid c$ แล้ว $a \mid(b x+c y)$ สำหรับทุกจำนวนเต็ม $x, y$
	\end{enumerate}
\end{theorembox}

